\section{Introduction}
\label{sec:intro}

Real-time smart city control loops, autonomous logistics, and ride-hailing platforms in the Moscow metropolitan agglomeration demand high-throughput routing engines capable of handling thousands of queries per second under continuous traffic fluctuations. 
Dynamic urban environments exhibit frequent congestion spikes and road closures. 
These rapid updates render static pathfinding acceleration techniques, such as Contraction Hierarchies, impractical, as the high cost of recomputing shortcut structures prevents real-time adaptation. 

We model the total routing execution time $T$ as the product of the search iteration count $N_{\text{iter}}$ and the average step latency $t_{\text{iter}}$:
\begin{equation}
T = N_{\text{iter}} \cdot t_{\text{iter}}
\label{eq:latency_model}
\end{equation}
The search algorithm (Dijkstra, A*, or ALT) determines the iteration count $N_{\text{iter}}$ by defining the relaxed search space envelope. 
The priority queue implementation determines the step latency $t_{\text{iter}}$ through memory layout, branch prediction efficiency, and cache hierarchy access patterns. 
Prior works focus on asymptotic complexity bounds. 
We focus on the interaction between algorithmic search space restriction and low-level hardware execution efficiency.

We evaluate seventy-seven configuration variants (seven priority queues across three algorithms with $1 + 2 \times 5$ weight parameters) on an Edge-Based Graph representation of the Moscow metropolitan network. 
Unlike the massive continental networks evaluated in general priority queue benchmarks, this subgraph focuses on representing high-density urban maneuvers. 
Our evaluation reveals a phenomenon we define as microarchitectural isomorphism: while priority queues scale the step latency constant $t_{\text{iter}}$, they exert zero influence on the macro-topological search space relaxation. 
The complexity exponents $\alpha$ remain parallel across different heaps.

We propose a three-step filtering pipeline to identify the Pareto-optimal routing configurations. 
First, we apply algorithmic filtering based on complexity exponents $\alpha$ to eliminate search paradigms with excessive exploration envelopes. 
Second, we evaluate heuristic weight configurations based on route accuracy bounds and priority queue monotonicity invariants. 
Third, we analyze microarchitectural efficiency (roughness and tail latency tail kurtosis) to select hardware-stable queue layouts. 
This filtering pipeline discards suboptimal combinations. 
Our work identifies Radix Heap as the optimal priority queue for monotone ALT routing, and the SIMD 8-ary heap for weighted A* search.

